Put $\mu(z)=h(z,z^{-1})$.  We need the antidiagonal cocharacter in the
common derived group, not merely in the Hodge group; the determinant gives a
short way to see this.  Since $\GGm$ is connected,
$\mu(\GGm)\subset G_H^0$.  By \cref{prop:hodge-vector},
\[
 \det\mu(z)=z^{(-2)\cdot2+0\cdot3+2\cdot2}=1.
\]
The determinant is trivial on the semisimple group $G^*$ and therefore
descends to a character
\[
 \overline{\det}:G_H^0/G^*\longrightarrow\GGm.
\]
The quotient $G_H^0/G^*$ is the image of $Z_H^0$.  On $Z_H^0$ the
determinant is $c\mapsto c^7$; since the faithful scalar action embeds
$Z_H^0$ into $\GGm$, the induced character $\overline{\det}$ has finite
kernel.  Hence the image of the connected group $\GGm$ under $\mu$ in
$G_H^0/G^*$ is both connected and finite, and is therefore trivial.
Consequently
\[
 \mu(\GGm)\subset G^*,
\]
the common Hodge and perverse derived group.  Its weights on $V$ are
\begin{equation}
 -2,0,2\quad\text{with multiplicities }2,3,2.
\label{eq:hodge-weights}
\end{equation}

Applying the fibre functor to \cref{prop:alt-support} gives
\begin{equation}
 \operatorname{Hom}_G(V,\Lambda^2V)\ne0.
\label{eq:hom-exterior}
\end{equation}

\begin{proposition}[Connected Tannaka group]\label[proposition]{prop:connected-g2}
There is an isomorphism of pairs
\[
 (G^0,V|_{G^0})\simeq(G_2,V_7).
\]
\end{proposition}

\begin{proof}
By \eqref{eq:orthogonal} and \eqref{eq:connected-irreducible}, the connected
centre of $G^0$ acts by orthogonal scalars and is therefore trivial.  Thus
$G^0$ is semisimple.  An irreducible representation of a product of
semisimple groups is an external tensor product.  Since $7$ is prime and the
action is faithful, $G^0$ has one simple factor.

We prefer to make the short highest-weight list explicit rather than leave
the classification implicit in the tables of
\cite{Dynkin1957,FultonHarris1991}.  After passage to the simply connected
cover, the simple modules of dimension seven are
\[
\begin{array}{c|c|c}
\text{type}&\text{highest weight}&\text{self-dual}\\ \hline
A_1&6\omega_1&\text{yes}\\
A_6&\omega_1,\ \omega_6&\text{no; the two are dual}\\
B_3&\omega_1&\text{yes}\\
G_2&\omega_1&\text{yes}.
\end{array}
\]
No other simple type has an irreducible module of dimension seven.  The
orthogonal self-duality of $V$ excludes type $A_6$.  In type $A_1$, the
central element $-I\in \mathrm{SL}_2$ acts trivially on $\Sym^6(\CC^2)$, so the
faithful image is $\mathrm{PGL}_2$; in type $B_3$, the seven-dimensional vector
representation has faithful image $\mathrm{SO}_7$.  Thus the remaining possibilities
are $\mathrm{PGL}_2$ on $\Sym^6(\CC^2)$, $\mathrm{SO}_7$ on its standard module, and $G_2$ on
$V_7$.  For $\mathrm{SO}_7$, the exterior
square of the standard module is the irreducible adjoint module, so
$\operatorname{Hom}_{\mathrm{SO}_7}(V,\Lambda^2V)=0$, contrary to
\eqref{eq:hom-exterior}.  For a nontrivial cocharacter of $\mathrm{PGL}_2$, the
weights on $\Sym^6(\CC^2)$ are seven distinct numbers.  For a primitive
cocharacter they are
\[
 -3,-2,-1,0,1,2,3,
\]
and every other nonzero cocharacter multiplies them by a nonzero integer.
The trivial cocharacter has one weight of multiplicity seven.  Neither
possibility agrees with \eqref{eq:hodge-weights}.  Hence $G^0=G_2$.
\end{proof}

\begin{corollary}[The full Tannaka group]\label[corollary]{cor:full-g2}
There is an isomorphism of pairs
\[
 (G,V)\simeq(G_2,V_7).
\]
In particular, $G$ is connected; this identifies the full Tannaka group, not
only its identity component.
\end{corollary}

\begin{proof}
By \cref{prop:connected-g2}, $G^0\simeq G_2$ on $V_7$.  Since
$G^0\mathrel{\triangleleft}G\subset \mathrm{O}(V)$,
\[
 G\subset N_{\mathrm{O}(V)}(G_2)=G_2\times\{\pm I\}.
\]
Indeed, conjugation gives a homomorphism from the normalizer to
$\operatorname{Aut}(G_2)$.  Since $\operatorname{Out}(G_2)=1$, multiplication
by a suitable element of $G_2$ puts every normalizing element in the
centralizer.  Schur's lemma makes that centralizer scalar, and its intersection
with $\mathrm{O}(V)$ is $\{\pm I\}$.  If $G\ne G^0$, choose
$g\in G\cap(-G_2)$ and write $g=(-I)h$ with $h\in G_2=G^0$.  Then
$-I=gh^{-1}\in G$.  But $-I$ acts by $-1$ on $V$ and by $+1$ on
$\Lambda^2V$, contradicting \eqref{eq:hom-exterior}.  Hence $G=G^0=G_2$.
\end{proof}

\begin{proof}[Proof of \cref{thm:main}]
Let $\cB^\circ$ be the nonempty open subset obtained in
\cref{thm:uniform-pair-splitting}.  The relative divisor $\cX\subset\cP$ is
symmetric, relatively ample, effective Cartier, flat, and has geometrically
integral fibres by the construction at the beginning of
\cref{sec:pair-splitting}.

Fix $b\in\cB^\circ(\CC)$.  The polarization type and the geometric
properties of $(P_b,X_b)$ follow from \cref{prop:theta-resolution} and the
construction of $L$.  The locus $\Omega_b$ in \cref{prop:hodge-vector} is
nonempty; for any $\alpha_b\in\Omega_b$, that proposition gives Hodge
multiplicities $(2,3,2)$.  The two pair components furnished by
\cref{thm:uniform-pair-splitting} imply, through \cref{prop:alt-support},
that $\IC_{X_b}$ is a direct summand of $\Alt^2(\IC_{X_b})$.  Finally,
\cref{cor:full-g2} identifies the full Tannaka group and tautological module
with $(G_2,V_7)$.  Since $b$ was arbitrary in $\cB^\circ(\CC)$, the theorem
follows.
\end{proof}

\appendix

\section{Exact norm identities and arithmetic specializations}\label{app:exact-norm}

This appendix contains the algebra used in \cref{sec:pair-splitting}.  We
keep the full residual octic out of the main text, but display every identity
that is used later: the normalized conic coefficients, the square-class
identity, and the two finite-field specializations.

\subsection*{A.1. Normalized conic coefficients}
On the normalized chart $e_1=v=1$ the shifted curve is
\begin{equation}
 Y^2+2cY
 =\lambda x^3+\bigl(1-\lambda(1-3e_2+3e_3)\bigr)x^2
 +\bigl(2c+\lambda(e_2^3-3e_2e_3+3e_3^2)\bigr)x
 -\lambda e_3^3.
\label{eq:appendix-shifted-curve}
\end{equation}
For the section $s=us_0+s_1+zs_2$ one has
\[
 \alpha=u(x-Y)+(Y-e_3)-zY,\qquad
 \beta=e_2+z(x-e_3),\qquad
 \gamma=-1+e_2z.
\]
Substitute these expressions into
$\operatorname{Nm}(s)=\alpha^3+\beta^3Y+\gamma^3Y^2-3\alpha\beta\gamma Y$,
reduce $Y^2$ using \eqref{eq:appendix-shifted-curve}, and divide by the fixed
line $Y-x$.  The quotient
$Q_s=q_{00}+q_{01}x+q_{02}Y+q_{11}x^2+q_{12}xY+q_{22}Y^2$ has
\begin{align*}
q_{00}&=\frac{2c+\lambda(e_2^3-3e_2e_3-3e_3^2u+3e_3^2)}{\lambda},\\
q_{01}&=\frac{1+\lambda(3e_2+3e_3u^2-3e_3-1)}{\lambda},\\
q_{11}&=-(u-1)(u^2+u+1),
\end{align*}
and
\begin{align*}
\lambda q_{02}={}&2cz^3+e_2^3\lambda z^3+3e_2^2\lambda uz-3e_2^2\lambda z
 -3e_2e_3\lambda uz^2-3e_2e_3\lambda z^3+3e_2e_3\lambda z^2\\
&-3e_2\lambda u+3e_2\lambda-3e_3\lambda u^2-3e_3\lambda uz
 +6e_3\lambda u+3e_3\lambda z-3e_3\lambda-\lambda+1,\\
\lambda q_{12}={}&3e_2\lambda uz^2+3e_2\lambda z^3-3e_2\lambda z^2
 +2\lambda u^3+3\lambda u^2z-3\lambda u^2-3\lambda uz\\
&-\lambda z^3+\lambda+z^3,\\
\lambda q_{22}={}&-\lambda u^3-3\lambda u^2z+3\lambda u^2-3\lambda uz^2
 +6\lambda uz-3\lambda u\\
&-\lambda z^3+3\lambda z^2-3\lambda z+\lambda+z^3.
\end{align*}
These six coefficients determine $M_s$ and hence $F_8$.  In particular, for
\[
 M_s=\begin{pmatrix}
 q_{00}&q_{01}/2&q_{02}/2\\
 q_{01}/2&q_{11}&q_{12}/2\\
 q_{02}/2&q_{12}/2&q_{22}
 \end{pmatrix}
\]
the only determinant expansion used below is
\begin{equation}
 4\det M_s=
 4q_{00}q_{11}q_{22}-q_{00}q_{12}^2-q_{11}q_{02}^2
 -q_{22}q_{01}^2+q_{01}q_{02}q_{12}.
\label{eq:appendix-determinant-expansion}
\end{equation}

